\documentclass[11pt]{amsart}
\usepackage{geometry}                % See geometry.pdf to learn the layout options. There are lots.
\geometry{letterpaper}                   % ... or a4paper or a5paper or ... 
%\geometry{landscape}                % Activate for for rotated page geometry
%\usepackage[parfill]{parskip}    % Activate to begin paragraphs with an empty line rather than an indent
\usepackage{graphicx}
\usepackage{amssymb}
\usepackage{epstopdf}
\usepackage{color}
\usepackage{colortbl}
\DeclareGraphicsRule{.tif}{png}{.png}{`convert #1 `dirname #1`/`basename #1 .tif`.png}

\title{PROBLEMA 4.14}
%\date{}                                           % Activate to display a given date or no date

\begin{document}
\maketitle
%\section{}
%\subsection{}

\noindent Se tienen los costos de producción de tres departamentos (dulces, bebidas y conservas), correspondientes a los 12 meses del año anterior. Construya un diagrama de flujo que pueda proporcionar la siguiente información:\\

\noindent \quad a) ¿En qué mes se registró el mayor costo de producción de dulces?\\

\noindent \quad b) Promedio anual de los costos de producción de bebidas.\\

\noindent \quad c) ¿En qué mes se registró el mayor costo de producción en bebidas, y en qué mes el menor costo?\\

\noindent \quad d) ¿Cuál fue el rubro que tuvo el menor costo de producción en diciembre?\\

\noindent Dato: PROD [1..12,1..3] (arreglo bidimensional de tipo real que almacena los costos de producción de tres departamentos en los 12 meses del año anterior).\\

\noindent \textbf {Explicación de las variables} \\

\noindent \quad I y J: Variables de tipo entero. Se usan como variables de control de varios ciclos y como índices de los arreglos.\\

\noindent \quad PROD: Arreglo bidimensional de tipo real..\\

\noindent \quad MAYDUL: Variable de tipo real. Se utiliza para almacenar el mayor costo de producción de dulces a lo largo del año.\\

\noindent \quad MES: Variable de tipo entero. Almacena el número de mes en el cual se registró el mayor costo de producción de dulces.\\

\noindent \quad SUM: Variable de tipo real. Se usa para sumar los costos mensuales de producción de bebidas, y para posteriormente calcular el promedio anual.\\

\noindent \quad MESMAY: Variable de tipo entero. Se utiliza para almacenar el mes en el cual se registró el mayor costo de producción de bebidas.\\

\noindent \quad MESMEN: Variable de tipo entero. Guarda el mes en el cual se registró el menor costo de producción de bebidas.\\

\noindent \quad MAYBEB: Variable de tipo real. Almacena el mayor costo mensual de producción de bebidas.\\

\noindent \quad MENBEB: Variable de tipo real. Almacena el menor costo mensual de producción de bebidas.\\

\noindent \quad MENOR: Variable de tipo real. Se usa para almacenar el menor costo de producción en el mes de diciembre.\\

\noindent \quad DEP: Variable de tipo entero. Almacena el número de departamento en el cual se registró el menor costo de producción en el mes de diciembre.\\

\end{document}  
